% The cell list in detail: (a) how the structure is built, (b) how a pair met
% in several cells is emitted once.
%
% No prose inside the drawing. What the panels mean belongs in the caption,
% where it can be read at caption size and does not compete with the marks.
%
% Grids are drawn line by line rather than with TikZ's `grid`, which places its
% lines at multiples of the step measured from the origin and not from the
% corner of the rectangle given to it.
\begin{tikzpicture}[
  font=\scriptsize,
  >={Stealth[length=1.6mm]},
  lbl/.style={font=\scriptsize, text=sccdInk, inner sep=1.5pt},
  tiny/.style={font=\tiny, text=sccdInk, inner sep=1pt},
  ttl/.style={font=\scriptsize\bfseries, text=sccdInk, inner sep=0pt},
  slot/.style={draw=sccdInk!45, line width=0.3pt, minimum width=3.6mm,
               minimum height=3.6mm, inner sep=0pt, font=\tiny, text=sccdInk},
  lead/.style={draw=sccdInk!55, line width=0.3pt},
]

% ============================== (a) construction =============================
% Three boxes over a 3x3 grid, and the two arrays they bin into. Cells are
% numbered row-major from the bottom left, which is the order cell_of gives.
\begin{scope}
  \node[ttl, anchor=south west] at (0,3.30) {(a) construction};

  \foreach \x in {0,0.70,1.40,2.10}
    \draw[draw=sccdInk!30, line width=0.3pt] (\x,0.90) -- (\x,3.00);
  \foreach \y in {0.90,1.60,2.30,3.00}
    \draw[draw=sccdInk!30, line width=0.3pt] (0,\y) -- (2.10,\y);
  \draw[draw=sccdInk!45, line width=0.4pt] (0,0.90) rectangle (2.10,3.00);

  \draw[draw=sccdTeal, line width=0.9pt] (0.15,1.05) rectangle (0.95,1.50);
  \draw[draw=sccdPlum, line width=0.9pt] (0.85,1.75) rectangle (1.55,2.20);
  \draw[draw=sccdAmber, line width=0.9pt] (1.50,2.42) rectangle (1.95,2.88);
  \node[lbl, text=sccdTeal, anchor=west] at (0.99,1.28) {$a$};
  \node[lbl, text=sccdPlum, anchor=west] at (1.59,1.98) {$b$};
  \node[lbl, text=sccdAmber, anchor=west] at (1.99,2.65) {$c$};

  % The cell numbers go on last, so a box cannot hide one.
  \foreach \c/\cx/\cy in {0/0.04/0.94, 1/0.74/0.94, 2/1.44/0.94,
                          3/0.04/1.64, 4/0.74/1.64, 5/1.44/1.64,
                          6/0.04/2.34, 7/0.74/2.34, 8/1.44/2.34}
    \node[tiny, text=sccdInk!55, anchor=south west, fill=white,
          inner sep=0.5pt] at (\cx,\cy) {\c};

  % --- the two arrays ---
  \node[lbl, anchor=east] at (-0.06,0.44) {\texttt{cellptr}};
  \foreach \i/\v in {0/0, 1/1, 2/2, 3/2, 4/2, 5/3, 6/4, 7/4, 8/4, 9/5}
    \node[slot, anchor=south west] at ({0.02 + \i*0.38},0.26) {\v};

  \node[lbl, anchor=east] at (-0.06,-0.16) {\texttt{cellidx}};
  \foreach \i/\v/\col in {0/$a$/sccdTeal, 1/$a$/sccdTeal, 2/$b$/sccdPlum,
                          3/$b$/sccdPlum, 4/$c$/sccdAmber}
    \node[slot, text=\col, anchor=south west] at ({0.02 + \i*0.38},-0.34) {\v};

  % Which cell each run of cellidx belongs to.
  \foreach \i/\c in {0/0, 1/1, 2/4, 3/5, 4/8}
    \node[tiny, text=sccdInk!55, anchor=north] at ({0.21 + \i*0.38},-0.40) {\c};
  \node[tiny, text=sccdInk!55, anchor=east] at (-0.06,-0.50) {cell};
\end{scope}

% ============================ (b) one pair, one cell =========================
\begin{scope}[xshift=57mm]
  \node[ttl, anchor=south west] at (0,3.30) {(b) one pair, one cell};

  \fill[sccdTeal!12] (0.60,1.05) rectangle (1.80,2.85);

  \foreach \x in {0,0.60,1.20,1.80,2.40,3.00}
    \draw[draw=sccdInk!30, line width=0.3pt] (\x,0.45) -- (\x,3.15);
  \foreach \y in {0.45,1.05,1.65,2.25,2.85}
    \draw[draw=sccdInk!30, line width=0.3pt] (0,\y) -- (3.00,\y);
  \draw[draw=sccdInk!45, line width=0.4pt] (0,0.45) rectangle (3.00,3.15);

  \draw[draw=sccdPlum!55, line width=0.5pt, densely dashed] (1.20,1.65) rectangle (2.40,2.85);

  \fill[sccdAmber!45] (1.30,1.80) rectangle (1.70,2.30);

  \draw[draw=sccdTeal, line width=0.9pt] (0.70,1.20) rectangle (1.70,2.30);
  \draw[draw=sccdPlum, line width=0.9pt] (1.30,1.80) rectangle (2.30,2.80);

  % The two cells the pair is met in go on top of the boxes, or the box edges
  % running beside them leave nothing of either mark to see.
  \draw[draw=sccdInk, line width=1.1pt] (1.20,1.65) rectangle (1.80,2.25);
  \draw[draw=sccdInk!45, line width=0.8pt] (1.32,2.37) -- (1.68,2.73);
  \draw[draw=sccdInk!45, line width=0.8pt] (1.68,2.37) -- (1.32,2.73);
  \fill[sccdInk] (1.30,1.80) circle (1.6pt);

  \node[lbl, text=sccdTeal, anchor=east] at (0.66,1.32) {first};
  \node[lbl, text=sccdPlum, anchor=west] at (2.34,2.70) {second};

  % --- key in one column, so no label can run into the next ---
  \fill[sccdAmber!45] (0.00,0.14) rectangle ++(0.22,0.16);
  \node[lbl, anchor=west] at (0.26,0.22) {overlap of the two boxes};
  \fill[sccdInk] (0.11,-0.08) circle (1.6pt);
  \node[lbl, anchor=west] at (0.26,-0.08) {its minimum corner};
  \draw[draw=sccdInk, line width=1.0pt] (0.00,-0.38) rectangle ++(0.22,0.16);
  \node[lbl, anchor=west] at (0.26,-0.30) {the cell that emits the pair};
  \draw[draw=sccdInk!40, line width=0.7pt] (0.02,-0.60) -- (0.20,-0.44);
  \draw[draw=sccdInk!40, line width=0.7pt] (0.20,-0.60) -- (0.02,-0.44);
  \node[lbl, anchor=west] at (0.26,-0.52) {met again here, and dropped};
\end{scope}
\end{tikzpicture}
